\documentclass[10pt,a4paper]{amsart}
\usepackage[margin=19mm]{geometry}
\usepackage{amsmath,amssymb,mathtools}
\usepackage[scr=rsfs,cal=euler]{mathalfa}
\usepackage{xfrac}
\usepackage[unicode,hidelinks,bookmarks=false]{hyperref}
\newcommand{\ie}{i.e.\ }
\newcommand{\cf}{cf.\ }
\newcommand{\resp}{respectively\ }
\newcommand{\Subsection}{\subsection}
\providecommand{\nobreakdash}{\nobreak\hbox{-}}
\setlength{\emergencystretch}{2em}
\multlinegap0pt
\usepackage{amssymb}
\usepackage{xcolor}
\usepackage{enumerate}
\newtheorem{lemma}{Lemma}
\newcommand{\Ed}{\mathbb{E}^{d}}
\newcommand{\aff}{\operatorname*{aff}}
\newcommand{\inte}{\operatorname*{int}}
\title{Characterizations of generalized convex bodies of revolution}
\author{M. Angeles Alfonseca, M. Cordier, E. Morales-Amaya, D. J. Verdusco Hern\'andez}
\date{Source: arXiv:2006.12646v5 (CC BY 4.0)}
\begin{document}
\maketitle

\noindent\emph{Source and licence.} Characterizations of generalized convex bodies of revolution. Version arXiv:2006.12646v5 (CC BY 4.0), distributed by arXiv under the Creative Commons Attribution 4.0 International licence, http://creativecommons.org/licenses/by/4.0/, as recorded in the arXiv licence metadata for this exact version and shown on the abstract page https://arxiv.org/abs/2006.12646v5. Full licence notice: \texttt{licenses/arxiv-2006.12646-CC-BY-4.0.txt}.\par\smallskip

\noindent\textit{Editorial scope.} Counted unit: the article's lemma kiara (source0.tex lines 458--461). The article's complete original proof of it is delivered with it (source0.tex lines 462--472, one \texttt{proof} environment of 11 lines). No mathematics is written, repaired or re-derived: the statement and the proof are the article's own lines, in the article's own notation, and no retained line is substituted or re-flowed.  Retained uncounted as the source-local context the proof needs: lines 329--347.  Nothing was omitted from any retained span, so the package carries no omission record.  No retained line contains a citation, so the wrapper carries no bibliography and no bibliography entry.  No retained line contains a figure environment, an image-inclusion command or drawing code, so no image is delivered and the package declares no asset dependency.  Editorial formatting only, in addition: an AMS \texttt{amsart} preamble that restates the article's own declarations and macros used by the retained lines, namely the environments \texttt{lemma} on separate counters, together with the standard packages those lines need.  The retained environments print in this document as Lemma 1 in order of appearance, whereas the article prints the same items with the numbering of its own full document.  The complete native source of the article as distributed in the arXiv e-print for this exact version is bundled unchanged as \texttt{sources/103874/source0.tex}.
\begin{lemma}\label{contento}


\begin{itemize}

$~$

\item[I.] Let $\Phi$ be a plane convex figure and let   $\{L_1,\ldots, L_n \}$ be the collection of all its lines 
of symmetry. Then,  $\{L_1,\ldots,L_n \}$ is an $n$-star of lines.
 
\item[II.] Let $K \subset \Ed$ be a convex body and let $\{H_i \}$ be a sequence of hyperplanes that intersect $\inte K$. Suppose that $H_i \rightarrow H$, $L_i \subset H_i$ is a $(d-2)$-plane of symmetry ($p_i \in H_i$ is a center of symmetry) of $H_i \cap K$, and $L_i \rightarrow  L$ ($p_i \rightarrow  p$); then, $L$ is a $(d-2)$-plane of symmetry ($p$ is center of symmetry) of $H \cap K$. 

\item[III.] Let $K\subset \Ed, d\geq3,$ be a convex body. Suppose that $\{L_n\} \subset \Ed$ is a sequence of axes (hyperplanes) of symmetry of $K$ and $L$ is a line (hyperplane) such that $L_n \rightarrow L$. Then $L$ is an axis (hyperplane) of symmetry of $K$.

\item[IV.] Let $\{K_i\}$ be a sequence of convex figures such that  $K_i \rightarrow K$ and, for every $i\in\mathbb{N}$, the figure $K_i$ has two lines of symmetry determining an angle $\theta_i$. If 
$\lim_{i\to\infty}\theta_i=0$, then  $K$ is a circle.
\end{itemize}

\end{lemma}

\begin{lemma}\label{kiara} 
Let $K\subset \Ed, d\geq 3$, be a convex body. Let $\Pi$ be a $k$-dimensional subspace, $2\leq k\leq d-1$, and let $L'\subset \Ed \backslash \Pi$ be a line through the origin $o$, such that  $L'\not=\Pi^{\perp}$. 
Suppose that every line $L\subset \Pi$ through $o$ is an axis of symmetry of $K$ and, furthermore, $L'$ is an axis of symmetry of the body $K$. Then every line through $o$ and contained in $\aff\{\Pi, L'\}$ is an axis of symmetry of $K$.
\end{lemma}

\begin{proof}
Let $v$ be a unit vector parallel to $L'$. For each $u\in \mathbb{S}^{k-1} \subset \Pi$, we denote by $\alpha(u)$ the angle between the vector $u$ and $v$. Let $\delta_1$ and $\delta_2$ be the minimum and the maximum of $\alpha (\cdot)$.  Since $L'\subset \Ed \backslash \Pi$ and $L'\not=\Pi^{\perp}$ it follows that $\delta_1>0$ and $\delta_2<\pi/2$. By virtue of the continuity of $\alpha(u)$ and the compactness of $\mathbb{S}^{k-1}$, we have that 
$\alpha (\mathbb{S}^{k-1})=[\delta_1,\delta_2]$. We denote by $F$ the subset of irrational multiples of $\pi$ in $[\delta_1,\delta_2]$ and by $E$ the set $\alpha^{-1}(F)$. Since $F$ is dense in $[\delta_1,\delta_2]$, the set $E$ is dense in $\mathbb{S}^{k-1}$. For each $u\in E$, we denote by $L(u)$ the line through the origin with direction $u$, and by $\Pi(u)$ the plane 
$\aff \{u,v\}$. Since $u\in E$, we have that $\alpha(u)\in F$ and, consequently, 
$\Sigma(L(u), L')$ is a set of axes of symmetry of $K$ in $\Pi(u)$ which is dense in 
$\Omega(L(u), L')$. Thus, by Lemma \ref{contento} [III], every line in $\Pi(u)$ through $o$, is an axis of symmetry of $K$. 


Let $w$ be a unit vector in span$\{\Pi,L'\}$. We will show that  there exists an axis of symmetry $L_w$ of $K$ such that $L_w=span\{w\}$. Let $e \in \mathbb{S}^{k-1}$ such that 
$span\{e\}=\Pi \cap span\{v,w\}$. Since $E$ is dense in $\mathbb{S}^{k-1}$, there exists a sequence $\{e_i\}\subset E$ such that $e_i \rightarrow e$, when $i \rightarrow \infty$. It follows that $\Pi(e_i) \rightarrow span\{v,w\}$. Hence, we can find a sequence of lines $\{\Delta_i\}$, $\Delta_i \subset \Pi(e_i)$, such that $\Delta_i \rightarrow L_w=span\{w\}\subset span\{v,w\}$. By the argument in the previous paragraph, every line $\Delta_i$ is an axis of symmetry of the body $K$. Therefore, by Lemma \ref{contento} [III], $L_w$ is an axis of symmetry of $K$.
\end{proof}

\end{document}
